% !TeX root = thesis.tex
%
% Requires \usepackage{longtable} in the preamble.

\chapter{Transparency on the Use of AI}
\label{ap2:ai-transparency}

This appendix reports the use of Artificial Intelligence tools in the preparation of this dissertation, in
accordance with the MECD \emph{Guidelines for the Use of AI Tools} (V1, 24 May 2026).

\section{Tools Used}

Three tools were used, with different roles. \textbf{Claude} (Anthropic) was the main tool, used
throughout the writing and implementation phases. \textbf{Gemini} (Google) and \textbf{ChatGPT} (OpenAI)
were used occasionally, for brainstorming and for summarising papers during the literature review; no text
produced by either was incorporated into this dissertation.

All tools were used interactively. Source code, result files and working text were supplied to Claude, and
its responses were reviewed, corrected or rejected by the author before use.

\section{Use by Section}

Table~\ref{tab:ai_by_section} reports the use of AI at the granularity of the section, as required by the
guidelines. Unless stated otherwise, the tool used was Claude.

\begin{longtable}{p{3.4cm}p{3.0cm}p{7.3cm}}
\caption{Use of AI tools by section.}
\label{tab:ai_by_section} \\
\hline
\textbf{Section} & \textbf{Type of use} & \textbf{Description} \\
\hline
\endfirsthead
\multicolumn{3}{l}{\small\emph{continued from previous page}} \\
\hline
\textbf{Section} & \textbf{Type of use} & \textbf{Description} \\
\hline
\endhead
\hline
\endfoot

Chapter~\ref{chap:sota}, all sections & Conceptual support, writing support & Gemini and ChatGPT were used
to summarise papers during the literature review. Claude was used to produce
Section~\ref{sec:tree_ensembles} and to rewrite Section~\ref{sec:training}, and for reformulation
elsewhere. \\

Chapter~\ref{chap:methodology}, all sections & Writing support, conceptual support & Reformulation for
clarity, and identification of passages where the text no longer matched the implemented code. The two
figures in Sections~\ref{sec:stl_models} and~\ref{sec:mtl_models} were generated from the author's
specification. \\

Chapter~\ref{chap:results}, Chapter~\ref{chap:discussion}, Chapter~\ref{chap:conclusions} & Content
generation & Text produced with the tool from the result files of the
author's own training runs and from the author's analysis of them. The author defined the structure and
the arguments, and revised the text over several iterations. \\

Sections~\ref{sec:motivation} and~\ref{sec:problem} & Content generation & Text produced with the tool from the
author's instructions on the motivation and on the problems to be addressed. \\

Abstract, Resumo, Acknowledgments & Content generation & Text produced with the tool from the completed chapters and
from a list of names supplied by the author. The Portuguese text was checked by the author, a native speaker. \\

Appendix~\ref{ap1:eda_plots}, Appendix~\ref{app:fault_plots}, this appendix & Content generation &
Text produced with the tool and reviewed by the author. \\

Supporting code & Code generation & Modifications to the fault injection framework, alignment of the two
training notebooks, and generation of the notebooks that produce the figures and LaTeX tables. \\
\end{longtable}

\section{Representative Prompts}

The interaction was conversational and iterative. The prompts below are representative of each category of
use, reproduced as issued, in the mixture of Portuguese and English in which the work was conducted.

\begin{itemize}
    \item \emph{``como é que funciona a correlação?''}
    \item \emph{``esta tabela esta correta?''}
    \item \emph{``what is wrong here?''}
    \item \emph{``quero simplificar a injeção correlacionada. quero que o propagation scale seja igual em
          todas as falhas e que seja de 0.9''}
    \item \emph{``podes gerar uma imagem para o Single-Task Deep Learning Models''}
    \item \emph{``quero simplificar muito mais este paragrafo para nao confundir o leitor''}
\end{itemize}

\section{Examples of Iteration}

\subsection*{Section~\ref{sec:problem}, Problem Definition}

The author specified that the section should present the absence of fault data as one problem, and that a
second problem was needed to justify the use of machine learning. The first version was rejected as too
elaborate; a simplified version was requested. The author then rejected the second paragraph, instructing
that it should justify deep learning through the complexity of the interaction between sensors rather than
through the inadequacy of fixed limits. The author edited the resulting text directly and requested a final
correction pass.

\subsection*{Chapter~\ref{chap:discussion}, Discussion}

The first version contained five sections. The author removed one section and several paragraphs and
requested a shorter chapter. The condensed version was rejected as still too elaborate, and a simpler
version was produced. The author then instructed that one specific argument be removed, and the affected
section was rewritten around a different argument supported by the same results.

\subsection*{Section~\ref{sec:training}, Training Deep Neural Networks}

The author asked for the section to be extended to introduce gradient descent, and for an assessment of
the accompanying figure. The first version included equations and identified three inconsistencies between
the figure and the text. The author requested a shorter version without equations, and the figure was
adjusted to match.

\section{Validation and Responsibility}

All numerical values reported in Chapter~\ref{chap:results} were obtained from the result files produced
by the author's own training runs and cross-checked against them. All code was executed by the author and
its outputs inspected. All bibliographic references suggested by the tools were verified against the
original publications, since AI tools can produce plausible but nonexistent references.

The author assumes full responsibility for all content presented in this dissertation, including text,
code, figures and tables, and is able to explain and justify every section, result and implementation
decision it contains.